\documentclass[11pt]{article}
\usepackage{amsmath,amssymb,amsthm,hyperref}
\newtheorem{proposition}{Proposition}
\begin{document}
\title{Complex Border Products and Admissible Real Factors}
\author{}
\date{November 23, 2026}
\maketitle

\paragraph{The construction inspected.}
Mrinal Kumar's \emph{On the power of border of depth-3 arithmetic
circuits}, Theorem 1.5 and its proof, starts with a homogeneous
polynomial $P=\sum_{i=1}^T L_i^d$ over $\mathbb C$ and forms
\[
 A_\epsilon(x)=\prod_{i=1}^T(1+\epsilon^dL_i(x)^d)
                   =1+\epsilon^dP(x)+O(\epsilon^{2d}).
\]
Each degree-$d$ factor splits over $\mathbb C$, so
$(A_\epsilon-1)/\epsilon^d$ is a difference of two products of
affine forms and tends coefficientwise to $P$. The nonconstant
product has $dT$ factors. Thus two product gates do not mean two
heads. More seriously, the factors produced by this splitting need
not be real or positive with the manuscript's slope orientation.

The following elementary invariant shows why replacing that
complex splitting by a product of admissible real denominators
cannot preserve general Boolean signs in this particular construction.

\begin{proposition}[A product threshold obstruction]
Let $D(x)$ be a finite product of positive affine functions on
$\{0,1\}^n$, each having all variable coefficients nonnegative
or all variable coefficients nonpositive. For a positive constant
$c$, neither $D-c$ nor $c-D$ can strictly sign represent parity
on $n\ge3$ variables, or its complement.
The assertion also holds for the stricter admissibility conditions
in the manuscript.
\end{proposition}

\begin{proof}
Fix any two coordinates and any assignment of the others. On this
square write one denominator's four values as
$a,a+b,a+d,a+b+d$. They are positive, and $bd\ge0$. Hence
\[
                  a(a+b+d)\le(a+b)(a+d).
\]
Multiplication over the factors gives
\[
                          D_{00}D_{11}\le D_{01}D_{10}.
\tag{1}
\]
In particular it is impossible that both $D_{00},D_{11}>c$
and both $D_{01},D_{10}<c$, since these inequalities would make
the left side of (1) exceed $c^2$ and the right side less than $c^2$.

For parity on at least three variables, fix all but two variables.
Changing one of the fixed bits switches which diagonal has the
positive label. Whichever overall output orientation is chosen,
one of these faces would demand exactly the forbidden pattern
for the predicate $D>c$. This proves the assertion.
\end{proof}

The same argument excludes a limiting construction using such
products if its limiting score has the required strict signs.
There are only finitely many cube vertices, so convergence to a
strict separator would preserve all signs for sufficiently small
positive parameter, contradicting the proposition. This requires
only convergence of the actual scores on the cube.

\paragraph{Scope.}
The result does not exclude a difference of two nonconstant
products, or a sum of affine fractions. Nor does it refute Kumar's
theorem, which is explicitly over the complex numbers and charges
formal degree separately. The gap is a real, admissible,
head-count-controlled transfer; no such transfer is asserted here.
The denominator product square inequality also follows from the
existing Lean results in
\texttt{Results/DenominatorProductSubmodularity.lean}.

\paragraph{A separate rational expansion.}
Neil Epstein's \emph{Rational Functions as Sums of Reciprocals of
Polynomials} uses successive polynomial division. Its main theorem
writes a proper univariate fraction as
\[
 \frac{f}{g}=\frac1{h_0}-\frac1{h_0h_1}
       +\cdots+\frac{(-1)^r}{h_0\cdots h_r},
 \qquad r\le\deg f,
\]
where the positive degrees of the $h_i$ increase strictly. The
denominators are products of polynomials of increasing degree;
they are not asserted to be affine, positive, or one oriented.
Charging each displayed summand as one manuscript head is therefore
invalid. The theorem supplies no target head bound.

\paragraph{Primary sources.}
\url{https://mrinalkr.bitbucket.io/papers/topfanin-formaldeg_revised.pdf}

\url{https://math.gmu.edu/~nepstei2/ratfuns-accepted.pdf}
\end{document}
